\section{Open Questions}
\label{sec:open-questions}

The following five questions were left open by the paper and by our replication. Each has a concrete, CPU-feasible next step using free tools (PySCF, Qiskit statevector or Aer noisy simulator, block2 DMRG).

\subsection*{Q1: Strongly-correlated fragments}
\textbf{Question:} Does LAS-UCC retain chemical accuracy when the fragments themselves are strongly correlated (multi-bond breaking within a fragment; transition-metal $d$-shells)?

\textbf{Basis:} The paper's benchmarks use weakly-correlated H$_2$/ethylene fragments. If per-fragment CASCI is not converged, on-top UCC cannot fix it.

\textbf{Next step:} LAS-UCC on Cr$_2$ (two Cr atoms as fragments, CAS(6,6) per Cr, 6-31G) across the 3.5~\AA{} avoided crossing. Compare to DMRG($m{=}1000$) and to LASSCF alone.

\subsection*{Q2: Composition with DMET / EDMFT}
\textbf{Question:} How does LAS-UCC compose with density-matrix embedding theory (DMET) or extended DMFT --- which already give fragment definitions and inter-fragment couplings but stop at mean-field bath?

\textbf{Basis:} Structurally similar, complementary --- LAS-UCC's on-top UCC could replace DMET's mean-field-bath approximation. Paper does not connect.

\textbf{Next step:} DMET-style bath-embedded UCC on the 1D Hubbard model at half-filling, $U/t \in \{2,4,8\}$. Compare per-fragment energies + double occupancies to DMET-CCSD and DMRG.

\subsection*{Q3: Real catalyst / drug-molecule application}
\textbf{Question:} Can LAS-UCC accelerate the CASSCF-quality bottleneck in a real catalyst or drug workflow, and on which specific molecule?

\textbf{Basis:} Paper demonstrates only on H$_4$ and butadiene. Real CASSCF-bottleneck cases (Fe-Mo cofactor, [4Fe-4S] clusters, methylcobalamin, heme O$_2$ binding) are untouched.

\textbf{Next step:} LAS-UCC on methylcobalamin Co-C bond homolysis with per-fragment CAS(6,6), compare BDE to the CASPT2 literature value (Kepp 2011, $\sim$35 kcal/mol) and CCSD(T)/def2-SVP.

\subsection*{Q4: Noise robustness of fragment recombination}
\textbf{Question:} How does the UCC-on-top recombination step degrade under realistic two-qubit-gate noise?

\textbf{Basis:} Paper's Fig.~5 gives only noiseless resource counts. UCCSD with 26+ parameters is noise-sensitive; noise-induced error could exceed the LASSCF error the method is trying to fix.

\textbf{Next step:} Rerun our (H$_2$)$_2$ VQE-UCCSD (\texttt{code/replicate\_las\_vqe.py}) under \texttt{qiskit-aer} density-matrix simulator + depolarizing two-qubit error $p \in \{10^{-4},10^{-3},5{\times}10^{-3},10^{-2}\}$, 8k shots. Identify the crossover $p$ above which LASSCF alone beats LAS-UCC.

\subsection*{Q5: Adaptive fragment-partitioning}
\textbf{Question:} How sensitive is LAS-UCC to the fragment-partitioning choice, and can partitions be chosen adaptively from the LASSCF 1-RDM instead of by chemical intuition?

\textbf{Basis:} Paper hand-picks fragments. For real molecules the choice is nontrivial and materially affects intra/inter-fragment correlation splitting. No partition-sensitivity study.

\textbf{Next step:} Take our (H$_2$)$_2$/6-31G data (\texttt{report/evidence/las\_6-31g.json}) and rerun the fragment-product surrogate under three partitions (2+2, 1+3, 4+0) at $r_{\text{inter}} \in \{0.6, 1.0, 1.5, 3.0\}$~\AA. Score partitions via inter-block mutual information of the LASSCF 1-RDM (Rissler-Schollw\"ock $s(1|2)$) and test whether min-MI = min-error.
